\section{Discussion}
\label{sec:results-overall-interpretation}
\subsection{What the controlled relevance comparison establishes}
For the implemented reference-cutoff procedures, reading the diagnostics together identifies weak discrimination and poor localization despite nominal-looking truth acceptance. This is the substantive meaning of the controlled relevance experiment. The compensating-noise construction keeps the documented treatment/endogeneity distribution fixed while changing the excluded signal. Consequently, the observed relevance patterns cannot be explained simply by a changing marginal spread of treatment. This control does not hold the entire observed sample fixed across settings, nor does it isolate the contribution of every computational feature. The experiment studies the complete retained procedures along this particular path.

The scientific value is therefore more specific than showing that weaker instruments make estimation harder. The same design permits apparently reassuring retention of the true candidate to be read alongside limited exclusion of false candidates and extensive acceptance on the maintained grid. These are complementary marginal summaries, not a joint probability, and the false-candidate rates are not size-adjusted comparisons. The controlled construction supports a cleaner description of their disagreement; it does not transform that description into a causal decomposition of nuisance estimation, weak identification and numerical error.

\subsection{Accuracy, retention and discrimination answer different questions}
Point accuracy evaluates the selected grid minimizer, whereas candidate-test behavior concerns the statistic at a specified coefficient. The two summaries use related data but different functionals of the computed profile. A procedure can select a comparatively accurate point while its statistic rejects the true candidate too frequently or retains many false candidates. The sample-size comparisons demonstrate why improving the first summary does not guarantee movement of truth acceptance toward the benchmark. They do not establish a convergence rate, and the available sample sizes cannot identify an asymptotic regime for these implementations.

Truth acceptance is equally incomplete as a description of information. It measures retention at the true effect under the specified reference rule, but does not say how much of the candidate domain is retained. The false-candidate summaries show discrimination at selected deviations, while accepted measure and full-grid frequency describe different aspects of the stored profile. Their combination distinguishes narrow retention around truth from diffuse acceptance across the computational domain. This distinction is useful even when validity is unresolved: the observed decisions remain describable, but a near-benchmark frequency is not promoted to a theorem about coverage.

The interpretation also depends on which inverted object is under discussion. Under a fully specified randomized candidate-test rule, truth acceptance can represent coverage of the corresponding inverted set. That event identity does not establish nominal or uniform validity. Here direct tests are evaluated separately from stored grid profiles, and fresh pivotal draws can differ between calls. Thus the reported direct-at-truth diagnostic is not automatically membership coverage of the separately stored finite-grid set. The Framework makes this distinction explicit; the integrated empirical comparison should not erase it.

\subsection{What the benchmark comparisons reveal}
Oracle-GMM is useful because it knows the designated relevant controls, yet it also loses accuracy and localization as relevance weakens. High-dimensional nuisance estimation is therefore not the sole explanation for deterioration. This comparison does not establish that the oracle implementation is ideal: it still uses the historical density proxy, matrix operations and candidate grid. Knowledge of the relevant controls removes a particular specification burden, not every inferential or computational difficulty. It is a benchmark within the experiment rather than a finite-sample lower bound on error or a certified test.

Likewise, the frequent RMSE advantage of preserved DML-IVQR over preserved Full-GMM does not identify the causal effect of regularization. Their nuisance dimensions, intercept treatment, residualization and fitting routes differ. Quantile-specific reversals remain part of the evidence, rather than exceptions to be discarded. The complete cell evidence is mapped in Appendix~\ref{app:evidence-map}. A controlled component comparison would require changing and documenting individual features in a separately specified study; that experiment has not been carried out here. No ranking of the retained implementations should be extended to all DGPs or all high-dimensional IVQR methods.

\subsection{Theoretical motivation and executable limitations}
Orthogonalization can reduce local sensitivity to nuisance errors under its conditions, but it cannot create the treatment--instrument information needed to distinguish structural candidates \parencite{chernozhukov_et_al_2018}. Identification-robust inference concerns validity under a procedure's own assumptions, not guaranteed discrimination \parencite{chernozhukov_hansen_2008,jun_2008}. These concepts explain the questions asked by the experiment. They do not certify that the preserved density proxy, projection and same-sample nuisance estimates satisfy the ideal score's conditions.

The retained mixed-density argument is especially important for that boundary. It concerns the residual density after integrating over treatment, not the density conditional on a realized positive treatment value and not the fitted-normal proxy. Its divergence blocks the displayed standard finite-density justification. It does not identify an alternative limiting distribution, prove that all possible inferential arguments fail, or make the recorded statistics mathematically meaningless. The implementation differences and selected numerical sensitivities were not varied independently to identify the source of a particular acceptance or rejection pattern. Missing cross-fitting alone is not such a diagnosis either.

Numerical certificates have a similarly specific role. The later investigations concern fixed objectives and exported inputs within selected cases. They document residual-rank and decision sensitivity and do not replace the historical Monte Carlo fits. A certificate for a specified problem is not a full-experiment error frequency, and improved objective accuracy need not imply identical residual signs or test decisions. Historical worker-specific penalty states and portable runtime reproduction also remain incomplete. Inspection of frozen evidence must therefore be distinguished from deterministic reporting, selected-case replay and exact historical stochastic replay.

\subsection{Finite domains and reporting implications}
Finite-grid accepted measure summarizes recorded binary decisions, rather than continuous boundaries. Full-grid acceptance is a different event, and expansion to a wider finite domain is a localization diagnostic rather than a search of the real line. Similar accepted measures can arise from different patterns of retained candidates; a nonaccepted boundary cannot rule out a disconnected component farther away. The preserved domain diagnostics are informative about these computational limits without establishing global boundedness, unboundedness or formal nonidentification.

For future Monte Carlo reporting, this study supports stating the decision rule, candidate domain and evaluation sequence alongside point-error summaries. It also motivates distinguishing raw false-candidate rejection from comparisons between tests of equal size. These are reporting implications drawn from this experiment, not a recommendation that a particular estimator or correction is universally preferable. Further comparisons would need their own DGPs, nuisance assumptions and randomization specifications. Within the present evidence, the defensible conclusion is that accuracy, truth retention and localization must be read together before a favourable summary is interpreted as informative inference.
